\documentclass{article}
\usepackage[utf8]{inputenc}
\usepackage{geometry}
\geometry{
	a4paper,
	total={170mm,257mm},
	left=20mm,
	top=20mm,
}
\usepackage{graphicx}
\usepackage{titling}
\usepackage[american,RPvoltages]{circuiTikz}
\usepackage{tikz, tikz-cd}
\usepackage{pgfplots}
\usepackage{siunitx}
\usepackage{float}
\usepackage{amsmath,amsfonts,stmaryrd,amssymb} % Math packages
\title{The Analysis and Design of Linear Circuits  }
\author{Rodrigo Castillejos}
\date{July 2024}

\usepackage{fancyhdr}
\fancypagestyle{plain}{%  the preset of fancyhdr 
	\fancyhf{} % clear all header and footer fields
	\fancyfoot[L]{\thedate}
	\fancyhead[L]{Problem Resolution}
	\fancyhead[R]{\theauthor}
}
\makeatletter
\def\@maketitle{%
	\newpage
	\null
	\vskip 1em%
	\begin{center}%
		\let \footnote \thanks
		{\LARGE \@title \par}%
		\vskip 1em%
		%{\large \@date}%
	\end{center}%
	\par
	\vskip 1em}
\makeatother

\usepackage{lipsum}  
\usepackage{cmbright}

\begin{document}
	
	\maketitle
	
	\noindent\begin{tabular}{@{}ll}
		Author & \theauthor\\
		Chapter &  Basic Circuit Analysis\\
		Problem & 2.68
	\end{tabular}
	
	\section*{Exercise 2-64}
	Ideally, a voltmeter has infinite internal resistance and can be placed across any device to read the voltage without affecting the result. A particular digital multimeter (DMM), a common laboratory tool, is connected across the circuit shown in figure 1. The expected voltage was \qty{10.2}{\volt}. However, the DMM reads \qty{7.61}{\volt}. The large, but finite, internal resistance of the DMM was "loading" the circuit and causing a wrong measurement to be made. Find the value of the internal resistance $R_M$ of this DMM.
 	
	\begin{figure}[H]
		\centering	
		\begin{tikzpicture}[scale=0.8]
			\draw[blue, fill=blue!5, dashed, thick, ] (7,-0.5) rectangle (10,5.5);
			\draw (0,0) to[V, l2=$V_{IN}$  and \qty{15}{\volt}] (0,5)
				  to[R, l=$R_1$, a=\qty{4.7}{\mega\ohm}, -*] ++(5,0) coordinate (p1)
				  to[R, l=$R_2$, a=\qty{10}{\mega\ohm}] ++(0,-5)
				  to[short] (0,0);
			\draw (p1) to[short] ++(4,0)
				  to[R, a=$R_M$] ++(0,-5)
				  to[short, -*] ++(-4,0);
			\draw (p1) node[above right, xshift=-20pt]{$V_{1} = \qty{7.61}{\volt}$};
		\end{tikzpicture}
		\caption{}	
		\label{fig:figura1}
	\end{figure}
	
	\section*{Answer}
	The unknown resistor $R_M$ is in parallel with $R_2$ resistor. Since voltages across parallel elements are equal, the voltage $V_1$ appears across both. We first define an equivalent resistance $R_{EQ} = R_M || R_2$ as
	
	\begin{align}
		R_{EQ} = \frac{R_2 R_M}{R_2 + R_M}
	\end{align}
	
	We write the voltage division rule in terms of $R_{EQ}$ as
	
	\begin{align}
		V_1 = \frac{R_{EQ}}{R_{EQ} + R_1} V_{IN} = \qty{7.61}{\volt}
	\end{align}
	
	Solving equation (2) for $R_{EQ}$ yields
	
	\begin{align}
		R_{EQ}V_{IN} = 	V_1(R_{EQ} + R_1) \nonumber \\
		R_{EQ}V_{IN} = R_{EQ}V_1 + R_1V_1 \nonumber \\
		(V_{IN} - V_1)R_{EQ} = R_1V_1 \nonumber \\
		R_{EQ} = \frac{R_1V_1}{V_{IN} - V_1} 
	\end{align}
	
	Therefore
	
	\begin{align}
		R_{EQ} = \qty{4.84}{\mega\ohm}
	\end{align}
	 
	Then, we sustitute the value presented in (4) on equation (1) ans solving for $R_2$
	
	\begin{align}
		R_2R_M = R_{EQ}(R_2 + R_M) \nonumber \\
		R_2R_M = R_2R_{EQ} + R_MR_{EQ} \nonumber \\
		(R_2 - R_{EQ})R_M = R_2R_{EQ} \nonumber
	\end{align}
	 
	 \begin{align}
	 	R_M = \frac{R_2R_{EQ}}{R_2 - R_{EQ}}
	 \end{align}

so

\begin{align}
	R_M = \frac{(\qty{10}{\mega\ohm}) (\qty{4.84}{\mega\ohm})}{\qty{10}{\mega\ohm} - \qty{4.84}{\mega\ohm}} = \qty{9.3798}{\mega\ohm}
\end{align}
	
\end{document}